%-------------------------------------------------------------------------------
% Chapter 6: Discussion and Contributions
%-------------------------------------------------------------------------------

\chapter{Discussion and Contributions}\label{chap:discussion}

This chapter synthesizes the empirical evidence established within this research to explicitly resolve the core research questions posed in the introduction. Rather than reiterating raw convergence performance metrics, this discussion interprets the underlying mathematical implications, architectural boundaries, and emergent behaviors observed throughout the non-deterministic automated algorithm design loops. By mapping our empirical observations directly back to the original research objectives, we validate the structural mechanism of noise-injected evolution. Furthermore, this chapter provides a white-box critical analysis detailing why our noise-evolved heuristics established superior cross-environment generalizability over deterministic baselines, addresses the economic limitations of commercial API-driven hyper-heuristics, and establishes the formal scientific contributions this thesis delivers to the field of autonomous metaheuristic discovery.

% ==============================================================================
% DISCUSSION & CONTRIBUTIONS ROADMAP (ADDRESSING THE RQ BOUNDARY LINK ISSUE)
% ==============================================================================
%
% % FUTURE SECTION 6.1: Direct Resolution of Research Questions
% % - CRITICAL DEFENSE NOTE: You must write explicit, dedicated sub-paragraphs here 
% %   re-stating each RQ and providing the summary answer based on Chapter 5 data.
% %
% %   - Link to RQ1 (The Noise Sandbox): Discuss the outcomes of your initial noise 
% %     distribution testing. Conclude whether Additive Gaussian Noise successfully 
% %     masked analytical gradients to establish a realistic physical noise floor, and 
% %     how the parameters chosen altered the evaluation loop difficulty.
% %
% %   - Link to RQ2 (Autonomous Synthesis Resilience): Answer if the LLM successfully 
% %     designed executable Python classes under stochastic pressure. Discuss the 
% %     nature of the emergent mathematical logic it wrote—specifically how it adapted 
% %     code statements to balance noise smoothing without getting permanently 
% %     stuck in rugged local valleys of the sub-sampled BBOB landscapes.
% %
% %   - Link to RQ3 (Cross-Environment Robustness): Resolve the ultimate validation. 
% %     Discuss the precise performance differential observed when the clean-trained 
% %     heuristics collapsed under unexpected noise anomalies, proving why noise-aware 
% %     evolution yields fundamentally superior generalizability metrics.
%
% % FUTURE SECTION 6.2: White-Box Interpretation of Emergent Algorithmic Logic
% % - Placeholder to analyze the computer science inside the AI's champion code.
% % - Discuss how the LLM responded to the evolutionary pressures of the stochastic loop.
% % - Address a fascinating aspect: Did the LLM merely memorize hardcoded human math 
% %   tricks (like moving averages), or did it combine code elements in unique structural 
% %   ways that challenge traditional human-designed noise-handling conventions?
%
% % FUTURE SECTION 6.3: Model Economics and Scalability Constraints
% % - Placeholder to address the resource bottleneck we tracked in our usage logs.
% % - Discuss token depletion trends, context wind footprint behaviors, and the 
% %   financial reality of scaling commercial API calls over multi-run parallel loops.
% % - Propose a formal transition framework to localized open-weights ecosystems 
% %   (e.g., Code Llama, StarCoder), detailing why private, quantized local processing 
% %   is required to make stochastic AAD loops industrially scalable.
%
% % FUTURE SECTION 6.4: Behavioral Diversity and the No Free Lunch Boundaries
% % - Placeholder to contextualize your results against the No Free Lunch (NFL) theorem.
% % - Discuss landscape tradeoffs: If a noise-evolved champion excelled on Group 4 
% %   multimodal structures, did it trade off performance efficiency on Group 1 separable 
% %   functions? Outline why tracking behavioural profiles outside a strict 1+1 loop 
% %   is necessary to manage early optimization convergence.
%
% % FUTURE SECTION 6.5: Comprehensive Summary of Contributions
% % - Final summary block stating the direct engineering assets left behind by this thesis.
% % - Explicitly list: (1) The validated stochastic IOH/BBOB software wrapper layer, 
% %   (2) The queryable local SQLite evolution archive capturing code transformation history, 
% %   and (3) The open-source suite of novel, robust metaheuristic Python scripts.
%
% ==============================================================================